%! Characteristics of Linear Functions — Unit 1, Lesson 1.0 — Experience & Formalize
%! (Activity · QuickNotes · Application · Check Your Understanding)
% Directory stays `experience/` (build identifier in shared/lesson.mk); the label the student
% and the teacher read is "Experience & Formalize".
% Math Medic sizing: 12pt body, OPEN white space after every question (no writing lines).
% \answerspace{H}{} reserves exactly H in the blank; the key passes the answer as the second
% argument, so the two files paginate identically by construction.
% Page budget (user, 2026-08-20): Activity 2pp max · QuickNotes 1/2pp · Application 1/2-1pp ·
% Check Your Understanding 1-2pp (practice, NOT scored — there is no homework).
\documentclass[12pt]{article}
\usepackage{algebra2-article}
\usepackage{algebra2-boxes}
\newcommand{\answerspace}[2]{\par\nopagebreak\noindent\begin{minipage}[t][#1][t]{\linewidth}%
  \color{keyred}\bfseries #2\end{minipage}\par}
\newcommand{\lgrid}[4]{%
  \draw[step=1, linegray!40, very thin] (#1,#3) grid (#2,#4);
  \draw[->, semithick, charcoal] (#1,0)--(#2,0) node[below right, font=\scriptsize]{$x$};
  \draw[->, semithick, charcoal] (0,#3)--(0,#4) node[left, font=\scriptsize]{$y$};}
% Temperature graph: hours h on the horizontal axis (0..7), temperature T on the vertical (-8..8).
% The h-axis is drawn at T=0 so the crossing is the point students circle.
\newcommand{\tempgrid}{%
  \draw[step=1, linegray!40, very thin] (0,-8) grid (7,8);
  \draw[->, semithick, charcoal] (0,0)--(7.6,0) node[right, font=\scriptsize]{$h$};
  \draw[->, semithick, charcoal] (0,-8.6)--(0,8.8) node[above, font=\scriptsize]{$T$ ($^\circ$C)};
  \foreach \x in {1,...,7} \node[charcoal, font=\tiny, below] at (\x,0){\x};
  \foreach \y in {-8,-6,-4,-2,2,4,6,8} \node[charcoal, font=\tiny, left] at (0,\y){\y};}

\begin{document}

\pageheader{Unit 1, Lesson 1.0}{Experience \& Formalize: Freezing Point}

\vspace{0.06in}

\begin{headlinebox}{forestbg}
\textbf{Freezing Point.} On a cold weekend the number that matters most is \emph{zero}: above it the
roads are wet, below it they are ice. Work the two temperature stories below with your group using
what you already know --- tables, rules, points --- and be ready to explain \emph{how you know} each
answer \textbf{from the graph}.
\end{headlinebox}

\vspace{0.04in}

% ══ Part 1 of 4 — ACTIVITY (20 min, groups of 3–4, prior knowledge only) ══════
% Two scenarios, same 2°/hr, same zero at h=3, opposite direction — the contrast IS
% the lesson. Both graphs are pre-drawn; students label and circle. Saturday = forest,
% Sunday = navy, matching the deck's launch and debrief frames.
\begin{scenariobox}[1. Saturday Morning]{forest}
At 6:00 a.m. on Saturday the temperature in Blacksburg was $-6\,^\circ$C. It rose steadily,
$2\,^\circ$C every hour, all morning (until noon). Let $h$ be the number of hours after 6:00 a.m.
and $T$ the temperature in $^\circ$C.

\begin{enumerate}[label=\textbf{\alph*.}, leftmargin=1.8em, itemsep=4pt]
  \item Complete the table.
    \begin{center}
    \renewcommand{\arraystretch}{1.5}
    \begin{tabular}{|c|c|c|c|c|c|c|c|}
    \hline
    $h$ & $0$ & $1$ & $2$ & $3$ & $4$ & $5$ & $6$ \\ \hline
    $T$ & $-6$ & $-4$ & \blank{1.0cm} & \blank{1.0cm} & \blank{1.0cm} & \blank{1.0cm} & \blank{1.0cm} \\ \hline
    \end{tabular}
    \end{center}
  \item Write a rule for the temperature $h$ hours after 6:00 a.m. \quad
        $T(h) = $ \blank{3.4cm}
\end{enumerate}

\begin{minipage}[t]{0.56\linewidth}
\vspace{0pt}
\begin{enumerate}[label=\textbf{\alph*.}, start=3, leftmargin=1.8em, itemsep=4pt]
  \item Circle where the line crosses the \emph{vertical} ($T$) axis:
        $(\blank{0.8cm},\ \blank{0.8cm})$. In the story, that point tells you:
        \answerspace{1.7cm}{}
  \item Circle where it crosses the \emph{horizontal} ($h$) axis:
        $(\blank{0.8cm},\ \blank{0.8cm})$. What time of day is that, and why is that moment
        special?
        \answerspace{1.7cm}{}
\end{enumerate}
\end{minipage}\hfill
\begin{minipage}[t]{0.40\linewidth}
\vspace{0pt}\centering
\begin{tikzpicture}[x=0.58cm, y=0.26cm]
  \tempgrid
  \draw[very thick, forest] (0,-6)--(7,8);
  \fill[charcoal] (0,-6) circle (2.8pt);
  \fill[charcoal] (3,0) circle (2.8pt);
\end{tikzpicture}
\end{minipage}

\begin{enumerate}[label=\textbf{\alph*.}, start=5, leftmargin=1.8em, itemsep=4pt]
  \item Every hour the temperature changes by \blank{1.2cm} $^\circ$C. Circle where that number
        shows up in your rule, then say how you \emph{see} it on the graph moving from one hour
        to the next.
        \answerspace{1.8cm}{}
  \item As the morning goes on, is the temperature going \textbf{up} or \textbf{down}? How does
        the \emph{graph} show that (not the story)?
        \answerspace{1.3cm}{}
        Below freezing: $T<0$ when \blank{2.6cm} \qquad
        Above freezing: $T>0$ when \blank{2.6cm}
\end{enumerate}
\end{scenariobox}

\boxguard[14]
\begin{scenariobox}[2. Sunday Evening]{navy}
At 4:00 p.m. on Sunday it was $6\,^\circ$C, and the temperature \emph{fell} steadily, $2\,^\circ$C
every hour, through the evening. Now $h$ counts the hours after 4:00 p.m.

\begin{minipage}[t]{0.54\linewidth}
\vspace{0pt}
\begin{enumerate}[label=\textbf{\alph*.}, leftmargin=1.8em, itemsep=5pt]
  \item Write a rule: \; $T(h) = $ \blank{3.0cm}
  \item The line crosses the $T$-axis at $(\blank{0.8cm},\ \blank{0.8cm})$ and the $h$-axis
        at $(\blank{0.8cm},\ \blank{0.8cm})$. What does each crossing tell you about Sunday
        evening?
        \answerspace{2.4cm}{}
\end{enumerate}
\end{minipage}\hfill
\begin{minipage}[t]{0.42\linewidth}
\vspace{0pt}\centering
\begin{tikzpicture}[x=0.58cm, y=0.26cm]
  \tempgrid
  \draw[very thick, navy] (0,6)--(7,-8);
  \fill[charcoal] (0,6) circle (2.8pt);
  \fill[charcoal] (3,0) circle (2.8pt);
\end{tikzpicture}
\end{minipage}

\begin{enumerate}[label=\textbf{\alph*.}, start=3, leftmargin=1.8em, itemsep=5pt]
  \item At 5:00 p.m. ($h=1$) the temperature is \blank{1.2cm} $^\circ$C. Is that number
        \textbf{positive} or \textbf{negative}? At that moment, is the temperature \textbf{rising}
        or \textbf{falling}? \\[4pt]
        So: can a temperature be \emph{positive} and \emph{falling} at the same time? Explain
        what each word is describing.
        \answerspace{2.2cm}{}
  \item Above freezing: $T>0$ when \blank{3.0cm} \qquad
        Below freezing: $T<0$ when \blank{2.6cm}
  \item \textbf{The story, then the math.} In the story, which hours actually make sense as
        inputs? \blank{4.0cm} \\[2pt]
        Now \emph{forget the story}: the rule $T(h) = 6 - 2h$ accepts \emph{any} number $h$ and its
        line runs forever both ways, so all the possible inputs are \blank{2.8cm} and all
        the possible outputs are \blank{2.8cm}. Why are those two answers different?
        \answerspace{1.5cm}{}
\end{enumerate}
\end{scenariobox}

% ══ Part 2 of 4 — QUICKNOTES (13 min debrief fills this; target half a page) ══
\boxguard[14]
\begin{tcolorbox}[enhanced, colback=sky, colframe=navy, boxrule=0.8pt, arc=2mm,
    left=3mm, right=3mm, top=1.5mm, bottom=1.5mm, breakable,
    fonttitle=\bfseries\small\color{white}, title={QuickNotes: Characteristics of a Linear Function},
    attach boxed title to top left={yshift=-2mm, xshift=4mm},
    boxed title style={colback=navy, colframe=navy, arc=1.5mm}]
\small
\begin{minipage}[t]{0.31\linewidth}
\vspace{2pt}\centering
\begin{tikzpicture}[scale=0.26]
  \lgrid{-2}{6}{-6}{6}
  \draw[very thick, forest] (-1,-6)--(5,6);
  \draw[semithick, navy, dashed] (2,0)--(3,0)--(3,2);
  \node[navy, font=\tiny] at (2.5,-0.7){run $1$};
  \node[navy, font=\tiny, right] at (3.05,1){rise $2$};
  \fill[charcoal] (0,-4) circle (3.4pt);
  \fill[charcoal] (2,0) circle (3.4pt);
\end{tikzpicture}\\[1pt]
{\footnotesize Example: $f(x) = 2x - 4$}
\end{minipage}\hfill
\begin{minipage}[t]{0.65\linewidth}
\vspace{0pt}
\begin{itemize}[leftmargin=1.1em, itemsep=3pt, topsep=0pt]
  \item A \textbf{linear function} has a \blank{2.2cm} rate of change, so its graph is a
        \blank{2.2cm}. Form: $f(x) = mx + b$.
  \item \textbf{Slope} $m = \dfrac{\Delta y}{\Delta x} = \dfrac{\text{rise}}{\text{run}}$. \;
        For $f$: $m = \blank{1.0cm}$.
  \item \textbf{$y$-intercept} $(0,b)$: set $x=0$. For $f$: $(0,\blank{0.9cm})$. \quad
        \textbf{$x$-intercept / zero} $(a,0)$: solve $f(x)=0$. For $f$: $(\blank{0.9cm},0)$.
\end{itemize}
\end{minipage}

\vspace{4pt}
\begin{itemize}[leftmargin=1.1em, itemsep=3pt, topsep=0pt]
  \item \textbf{Increasing / decreasing} --- read left to right: \;
        \mbox{$m>0$: \blank{2.0cm}} \quad \mbox{$m<0$: \blank{2.0cm}} \quad
        \mbox{$m=0$: \blank{2.0cm}}
  \item \textbf{Positive / negative:} $f(x)>0$ where the graph is \blank{1.6cm} the $x$-axis;
        $f(x)<0$ where it is \blank{1.6cm}. The sign can only switch at a \blank{1.6cm}.
        For $f$: positive when $x \blank{1.2cm}$, negative when $x \blank{1.2cm}$.
        \emph{Where it sits $\ne$ which way it moves} --- Sunday's line was positive \emph{and} falling.
  \item \textbf{Domain / range} of a line that is not horizontal: \blank{2.4cm} /
        \blank{2.4cm}. A horizontal line $f(x)=b$ has range \blank{1.4cm} and no
        zero (unless $b=0$). In a \emph{story}, the domain is only the inputs that make sense.
\end{itemize}
\end{tcolorbox}

% ══ Part 3 of 4 — APPLICATION (7 min, worked together; the modeling standard) ══
\boxguard[14]
\begin{notesbox}{Application: The Car-Wash Card}
We will do this one together. A car-wash gift card starts with \$40, and each wash costs \$8, so
the balance after $w$ washes is $B(w) = 40 - 8w$.

\begin{enumerate}[leftmargin=1.7em, itemsep=5pt]
  \item What does the slope $-8$ tell you about the card? And the $y$-intercept $(0,40)$?
    \answerspace{1.4cm}{}
  \item Find the zero of $B$ by solving $B(w) = 0$.
    \begin{work}
      40 - 8w &= 0 \\
          -8w &= -40 \\
            w &= 5
    \end{work}
    In the story, that zero means:
    \answerspace{1.4cm}{}
  \item Which inputs $w$ make sense? \blank{3.4cm} \;
        Is $B$ increasing or decreasing there? \blank{2.6cm}
  \item Suppose each wash cost \$10 instead of \$8. Would the card run out \emph{sooner} or
        \emph{later}? \blank{2.2cm} \; And on the graph, what changes --- where the
        line \emph{starts}, how \emph{steep} it is, or both? Say why.
    \answerspace{1.4cm}{}
\end{enumerate}
\end{notesbox}

% ══ Part 4 of 4 — CHECK YOUR UNDERSTANDING (10 min, practice, NOT scored) ═════
% This is the lesson's only practice set: there is no homework, so the six items
% below carry the whole A2.F.2 spread — graph, contrast pair, table, the constant
% special case with its boundary, a model in a fresh context, and the formative MC.
\boxguard[14]
\begin{notesbox}{Check Your Understanding \quad {\normalfont\itshape (practice --- not scored)}}
Six exercises --- the lesson's practice, and there is \textbf{no homework}. Work 1--3 with a partner,
then 4--6 on your own. For every answer, be ready to say \emph{how you know} it.

\begin{enumerate}[leftmargin=1.7em, itemsep=4pt]
  \item \textbf{Read the rule.} Complete the feature list for $f(x) = 2x + 2$.
    \begin{center}
    \renewcommand{\arraystretch}{1.4}
    \begin{tabular}{@{}l l l@{}}
      Slope: \blank{1.5cm} & $y$-intercept: \blank{1.8cm} & Zero: \blank{1.6cm} \\
      Domain: \blank{1.8cm} & Range: \blank{1.8cm} & Incr.\ or decr.? \blank{2.0cm} \\
      Positive when $x$ \blank{1.4cm} & Negative when $x$ \blank{1.4cm} & Sign flips at: \blank{1.6cm}
    \end{tabular}
    \end{center}

  \boxguard[12]
  \item \textbf{Same zero, opposite behavior.} Both lines have their zero at $x = 1$.
    \begin{center}
    \begin{minipage}[c]{0.34\linewidth}\centering
    \begin{tikzpicture}[scale=0.36]
      \lgrid{-3}{5}{-4}{4}
      \draw[very thick, forest] (-2.5,-3.5)--(4.5,3.5) node[above left, font=\scriptsize, forest]{$f$};
      \draw[very thick, navy] (-2.5,3.5)--(4.5,-3.5) node[below left, font=\scriptsize, navy]{$g$};
      \fill[charcoal] (1,0) circle (3.4pt);
    \end{tikzpicture}
    \end{minipage}\hfill
    \begin{minipage}[c]{0.62\linewidth}
    \begin{enumerate}[label=(\alph*), leftmargin=1.9em, itemsep=4pt]
      \item Increasing: \blank{1.2cm} \; Decreasing: \blank{1.2cm}
      \item $f(x)>0$ when $x$ \blank{1.4cm} \quad $g(x)>0$ when $x$ \blank{1.4cm}
      \item At $x=3$: $f$ is \blank{1.8cm} \; $g$ is \blank{1.8cm}
    \end{enumerate}
    \end{minipage}
    \end{center}
    (d) Why can a line's sign change at its \emph{zero} and nowhere else?
    \answerspace{1.4cm}{}

  \item \textbf{From a table.} A linear function is given below. \emph{Watch the step in $x$.}
    \begin{center}
    \renewcommand{\arraystretch}{1.3}
    \begin{tabular}{|c|c|c|c|c|}
    \hline
    $x$ & $0$ & $2$ & $4$ & $6$ \\ \hline
    $f(x)$ & $-6$ & $-3$ & $0$ & $3$ \\ \hline
    \end{tabular}
    \end{center}
    Slope $m = \dfrac{\Delta y}{\Delta x} = \blank{1.4cm}$ \quad
    $y$-intercept: \blank{1.8cm} \quad Zero: \blank{1.6cm} \\[4pt]
    Increasing or decreasing? \blank{2.4cm} \quad $f(x)<0$ when $x$ \blank{1.4cm}

  \boxguard[12]
  \item \textbf{The flat one.} The graph of $c(x) = -2$ is a horizontal line through $(0,-2)$.
    \begin{enumerate}[label=(\alph*), leftmargin=1.9em, itemsep=4pt]
      \item Slope: \blank{1.2cm} \; Incr., decr., or constant? \blank{2.0cm} \;
            Range: \blank{1.6cm} \; A zero? \blank{1.2cm}
      \item Is $c(x)$ positive or negative? \blank{3.0cm}
      \item For which value of $b$ does $c(x) = b$ \emph{have} a zero? \blank{1.8cm} \;
            And then how many zeros does it have? \blank{4.0cm}
    \end{enumerate}

  \boxguard[12]
  \item \textbf{Model it.} A pool already holds $600$ gallons when a hose is turned on; the hose
        adds $150$ gallons a minute, so $W(t) = 600 + 150t$ after $t$ minutes.
    \begin{enumerate}[label=(\alph*), leftmargin=1.9em, itemsep=4pt]
      \item Slope $150$ means \blank{4.4cm} \;
            $(0,600)$ means \blank{3.6cm}
      \item Find the zero of $W$ by solving $W(t) = 0$.
        \begin{work}
          600 + 150t &= 0 \\
                150t &= -600 \\
                   t &= -4
        \end{work}
        Does that zero mean anything in this story? Explain.
        \answerspace{1.2cm}{}
      \item Which inputs $t$ make sense? \blank{2.6cm} \; Is $W$ increasing or
            decreasing on them? \blank{2.4cm}
    \end{enumerate}

  \item \textbf{Multiple choice.} For the linear function $y = -3x + 6$, which statement is
        \emph{true}?
    \begin{enumerate}[label=\Alph*., leftmargin=2.2em, itemsep=1pt]
      \item[\textbf{A.}] The function is increasing. \hfill
            \textbf{B.} It is positive for all $x > 2$.
      \item[\textbf{C.}] The zero is at $x = 2$. \hfill
            \textbf{D.} The $y$-intercept is $(0,-3)$.
    \end{enumerate}
    Name one of the other three and say why it is false.
    \answerspace{1.4cm}{}
\end{enumerate}
\end{notesbox}

\boxguard[6]
\begin{spiralbox}
\textbf{Coming next --- Lesson 1.1: Linear Functions.} Today you \emph{read} a line's characteristics
off its graph, table, and rule. Next you run it backwards: slope and a point \emph{write} the
equation three ways, and shifting the parent line $y=x$ graphs it.
\end{spiralbox}

\end{document}
